\documentclass[border=8pt]{standalone}
\usepackage{ctex}
\usepackage{amssymb}
\usepackage{tikz}
\usetikzlibrary{arrows.meta}

% 画一个区间：L 形折线（∩ 形），#5=1 时端点空心，#5=0 时实心
\newcommand{\interval}[5]{%
  \draw[#4, line width=1.6pt]
    (#1,0) -- (#1,#3) -- (#2,#3) -- (#2,0);
  \ifnum#5=1
    \draw[white, line width=2.2pt] (#1,0) circle (2.4pt);
    \draw[#4, line width=0.9pt] (#1,0) circle (2.4pt);
    \draw[white, line width=2.2pt] (#2,0) circle (2.4pt);
    \draw[#4, line width=0.9pt] (#2,0) circle (2.4pt);
  \else
    \fill[#4] (#1,0) circle (2.4pt);
    \fill[#4] (#2,0) circle (2.4pt);
  \fi
}

\begin{document}
\begin{tikzpicture}[xscale=1.05, yscale=1.0, line cap=round, line join=round]

% ============ 第 0 步：端点的两种画法 ============
\begin{scope}[shift={(0,0)}]
  \draw[->, gray] (-3.8,0) -- (3.8,0);
  \foreach \x in {-3,-2,-1,1,2,3} \draw[gray] (\x,0.07) -- (\x,-0.07);
  \draw[white, line width=2.2pt] (-1.8,0) circle (2.6pt);
  \draw[black, line width=0.9pt] (-1.8,0) circle (2.6pt);
  \node[below=4pt, font=\footnotesize] at (-1.8,0) {空心：不含};
  \fill[black] (1.8,0) circle (2.6pt);
  \node[below=4pt, font=\footnotesize] at (1.8,0) {实心：含有};
  \node[above=4pt, font=\small] at (0,0.7) {\textbf{第 0 步：先认端点}};
\end{scope}

% ============ 第 1 步：A = B ============
\begin{scope}[shift={(0,-1.9)}]
  \draw[->, gray] (-3.8,0) -- (3.8,0);
  \foreach \x in {-3,-2,-1,1,2,3} \draw[gray] (\x,0.07) -- (\x,-0.07);
  \interval{-1.5}{1.5}{0.55}{blue}{1}
  \node[blue, above=3pt, font=\normalsize] at (-1.5,0.55) {$A$};
  \node[blue, above=3pt, font=\normalsize] at (1.5,0.55) {$B$};
  \node[below=2pt, font=\small] at (0,-0.85) {$A=B$：$p$ 是 $q$ 的\textbf{充要条件}};
\end{scope}

% ============ 第 2 步：A 真包含于 B ============
\begin{scope}[shift={(0,-3.8)}]
  \draw[->, gray] (-3.8,0) -- (3.8,0);
  \foreach \x in {-3,-2,-1,1,2,3} \draw[gray] (\x,0.07) -- (\x,-0.07);
  \interval{-2.5}{2.5}{1.05}{red!70}{1}
  \node[red!70, above=3pt, font=\normalsize] at (-2.5,1.05) {$B$};
  \interval{-0.5}{1.5}{0.55}{blue}{1}
  \node[blue, above=3pt, font=\normalsize] at (-0.5,0.55) {$A$};
  \node[below=2pt, font=\small] at (0,-0.85) {$A\subsetneqq B$：$p$ 是 $q$ 的\textbf{充分不必要条件}};
\end{scope}

% ============ 第 3 步：B 真包含于 A ============
\begin{scope}[shift={(0,-5.7)}]
  \draw[->, gray] (-3.8,0) -- (3.8,0);
  \foreach \x in {-3,-2,-1,1,2,3} \draw[gray] (\x,0.07) -- (\x,-0.07);
  \interval{-2.5}{2.5}{1.05}{blue}{1}
  \node[blue, above=3pt, font=\normalsize] at (-2.5,1.05) {$A$};
  \interval{-0.5}{1.5}{0.55}{red!70}{1}
  \node[red!70, above=3pt, font=\normalsize] at (-0.5,0.55) {$B$};
  \node[below=2pt, font=\small] at (0,-0.85) {$B\subsetneqq A$：$p$ 是 $q$ 的\textbf{必要不充分条件}};
\end{scope}

% ============ 第 4 步：互不包含 ============
\begin{scope}[shift={(0,-7.6)}]
  \draw[->, gray] (-3.8,0) -- (3.8,0);
  \foreach \x in {-3,-2,-1,1,2,3} \draw[gray] (\x,0.07) -- (\x,-0.07);
  \interval{-2.6}{0.4}{0.55}{blue}{1}
  \node[blue, above=3pt, font=\normalsize] at (-2.6,0.55) {$A$};
  \interval{-0.4}{2.4}{1.05}{red!70}{1}
  \node[red!70, above=3pt, font=\normalsize] at (-0.4,1.05) {$B$};
  \node[below=2pt, font=\small] at (0,-0.85) {互不包含：$p$ 是 $q$ 的\textbf{既不充分也不必要条件}};
\end{scope}

\end{tikzpicture}
\end{document}
